\chapter{Kernel Security}
\label{ch:kernel}

\epigraph{``The kernel is not a trust boundary. It is the trust anchor.''}{}

\learninggoal{
	\item Describe the privilege model of modern operating systems (rings, modes).
	\item Understand the kernel attack surface: system calls, device drivers, ioctl.
	\item Classify kernel vulnerability types and their exploitation primitives.
	\item Explain kernel defences: KASLR, SMEP, SMAP, KPTI, CFI.
	\item Trace the exploitation pipeline for a typical kernel vulnerability.
}

\section{Privilege Rings}

Processors enforce privilege through \emph{protection rings}. x86-64 provides four rings (0--3), though most operating systems use only two:

\begin{itemize}
	\item \textbf{Ring 0 (Kernel mode):} full access to all memory, I/O ports, and privileged instructions. The operating system kernel and device drivers run here.
	\item \textbf{Ring 3 (User mode):} restricted access. User-mode processes cannot execute privileged instructions (like \texttt{lgdt}, \texttt{lidt}), cannot access kernel memory, and cannot mask interrupts.
\end{itemize}

\begin{definition}[Privilege Escalation]
	\textbf{Privilege escalation} is the act of gaining higher privileges than those initially granted. In the context of kernel security, it typically means going from Ring 3 (user) to Ring 0 (kernel) code execution.
\end{definition}

The transition from user to kernel mode occurs only through controlled entry points: system calls, interrupts, and exceptions. The processor switches the privilege level and jumps to a handler whose address is stored in the Interrupt Descriptor Table (IDT) or via the \texttt{syscall} instruction (which reads from Model-Specific Registers).

\section{The Kernel Attack Surface}

A kernel vulnerability is a bug in kernel code that can be triggered from user space. The attack surface includes:

\subsection{System Calls}

Every system call is an entry point from user to kernel space. Linux has over 400 syscalls; each is a potential attack vector. A bug in any syscall implementation (e.g., a missing bounds check in \texttt{ioctl}) is a kernel vulnerability.

\subsection{Device Drivers}

Device drivers account for the majority of kernel code and a disproportionate share of vulnerabilities. Drivers are often written by third-party vendors with less rigorous review, and they operate with full kernel privileges.

\begin{lstlisting}[language=C, caption=Simplified vulnerable driver ioctl]
static long driver_ioctl(struct file *file, unsigned int cmd,
                         unsigned long arg) {
    char buf[64];
    switch (cmd) {
        case CMD_WRITE:
            if (copy_from_user(buf, (void __user *)arg, sizeof(buf)))
                return -EFAULT;
            // ... use buf ...
            break;
        case CMD_WRITE_LARGE:
            // BUG: no size check!
            copy_from_user(buf, (void __user *)arg, arg_size);
            break;
    }
}
\end{lstlisting}

The \texttt{CMD\_WRITE\_LARGE} path overflows the kernel stack buffer---a classic kernel stack buffer overflow.

\subsection{procfs and sysfs}

Pseudo-filesystems \texttt{/proc} and \texttt{/sys} expose kernel-internal data to user space. Reading or writing specific files invokes kernel callbacks. Bugs in these callbacks (race conditions, missing validation) are common.

\subsection{ioctl}

The \texttt{ioctl} system call on device files is notoriously complex. Each driver defines its own commands, argument structures, and validation logic---a rich source of vulnerabilities.

\section{Kernel Vulnerability Classes}

\subsection{Kernel Stack Overflow}

Analogous to user-space stack overflows but more severe: overflowing the kernel stack corrupts kernel control data.

\begin{lstlisting}[language=C, caption=Kernel stack overflow]
asmlinkage long sys_buggy_ioctl(unsigned int fd, unsigned int cmd,
                                 unsigned long arg) {
    char buf[64];
    // arg_size can be > 64!
    if (copy_from_user(buf, (void __user *)arg, arg_size))
        return -EFAULT;
}
\end{lstlisting}

Kernel stack is typically small (8--16 KB) and cannot grow. Stack overflows in the kernel are especially dangerous.

\subsection{Kernel Heap Overflow}

The kernel has its own allocators (slab, slub, slob). Heap overflows corrupt kernel objects:

\begin{lstlisting}[language=C, caption=Kernel heap overflow]
struct cred *cred = prepare_creds();
// cred is a kernel object allocated from slab
// Overflowing an adjacent slab object can overwrite cred->uid
\end{lstlisting}

\subsection{Use-After-Free in Kernel}

Kernel UAF occurs when a kernel object is freed but a reference remains:

\begin{lstlisting}[language=C, caption=Kernel use-after-free]
struct file *file = fget(fd);
// ... file might be closed and freed by another thread ...
fput(file);
// UAF: file->f_op->write(file, ...)
\end{lstlisting}

The classic kernel UAF pattern: race between two threads where one frees an object and the other accesses it.

\subsection{Double Fetch}

A double fetch occurs when kernel code reads user-space data twice, and the data changes between reads (a TOCTOU race):

\begin{lstlisting}[language=C, caption=Double fetch vulnerability]
long fixup_ioctl(void __user *arg) {
    struct user_data data;
    copy_from_user(&data, arg, sizeof(data));  // First fetch
    int size = data.size;
    // ... validation passes ...
    copy_from_user(kbuf, arg->buf, data.size); // Second fetch!
    // data.size may have changed!
}
\end{lstlisting}

A malicious thread modifies \texttt{data.size} between the two \texttt{copy\_from\_user} calls, bypassing the validation.

\section{Kernel Exploitation Primitives}

Kernel exploits aim to elevate privileges (get root) rather than spawn a shell. Common primitives:

\begin{longtable}{p{3.5cm}p{5cm}p{4cm}}
	\toprule
	\textbf{Primitive}                               & \textbf{Mechanism}                               & \textbf{Goal}         \\
	\midrule
	\texttt{commit\_creds(prepare\_kernel\_cred(0))} & Call these kernel functions                      & Root credentials      \\
	Modify \texttt{cred} struct                      & Overwrite \texttt{uid=0, gid=0} in process creds & Root                  \\
	Modify \texttt{task\_struct}                     & Overwrite \texttt{real\_cred} pointer            & Process creds         \\
	MMU manipulation                                 & Modify page tables                               & Code execution        \\
	modprobe\_path                                   & Overwrite path to execute as root                & Kernel-mode code exec \\
	\bottomrule
\end{longtable}

\subsection{commit\_creds Exploit}

The classic kernel exploit pattern:

\begin{enumerate}
	\item Gain arbitrary kernel memory write (via heap overflow, UAF, etc.).
	\item Write the address of \texttt{commit\_creds} and \texttt{prepare\_kernel\_cred(0)} into the kernel's function pointer table or overwrite a process's \texttt{cred} structure.
	\item Trigger execution flow that calls these functions.
	\item The current process now has root credentials.
\end{enumerate}

\subsection{Modprobe\_Path Technique}

A simpler technique that requires only an arbitrary write:

\begin{enumerate}
	\item Write to \texttt{/proc/sys/kernel/modprobe\_path} (or directly to kernel memory at the \texttt{modprobe\_path} symbol) to point to a script you control.
	\item Execute an unknown binary format, causing the kernel to invoke \texttt{modprobe\_path} as root.
	\item Your script runs as root.
\end{enumerate}

\section{Kernel Defences}

\subsection{KASLR}

Kernel Address Space Layout Randomisation (KASLR) randomises the base address where the kernel image is loaded. On x86-64, the kernel text is typically mapped at a random offset within a large region.

\subsection{SMEP and SMAP}

\begin{definition}[SMEP and SMAP]
	\textbf{SMEP} (Supervisor Mode Execution Prevention) prevents the kernel from executing code in user-space pages. \textbf{SMAP} (Supervisor Mode Access Prevention) prevents the kernel from accessing user-space data (except through explicit copy\_from/to\_user).
\end{definition}

SMEP forces kernel exploits to use ROP within kernel code (cannot redirect to user-space shellcode). SMAP prevents direct dereferencing of user-space pointers from kernel mode.

\subsection{KPTI}

Kernel Page Table Isolation (KPTI) separates kernel and user page tables. In user mode, only a minimal set of kernel pages are mapped (just enough to handle syscalls and interrupts). This mitigates Meltdown-style side-channel attacks and reduces the kernel's attack surface.

\subsection{Kernel CFI}

Clang's CFI can be applied to the kernel (CFI\_CLANG in Linux). Forward-edge CFI for the kernel restricts indirect call targets to valid functions, preventing many control-flow hijack exploits.

\subsection{seccomp for the Kernel}

While seccomp is a user-space defence, it also protects the kernel by reducing the syscall surface that attackers can use. A restricted seccomp profile means fewer kernel code paths are reachable from a compromised process.

\section{The Kernel Exploit Pipeline}

Modern kernel exploitation follows a pattern:

\begin{enumerate}
	\item \textbf{Gather information:} probe KASLR offset via side channels or uninitialised memory leaks.
	\item \textbf{Prepare user space:} mmap user pages at predictable addresses for ROP chains or data structures.
	\item \textbf{Trigger vulnerability:} invoke the buggy syscall or ioctl to gain a write primitive.
	\item \textbf{Bypass SMAP/SMEP:} use kernel ROP (gadgets within kernel code) or overwrite function pointers.
	\item \textbf{Escalate privileges:} call \texttt{commit\_creds(prepare\_kernel\_cred(0))} via ROP or function pointer hijack.
	\item \textbf{Restore execution:} return to user mode with root privileges.
\end{enumerate}

\begin{principle}[User-to-Kernel-to-User]
	Exploitation does not stop at kernel code execution. The exploit must return control to user space with the privilege elevation applied. This requires saving and restoring user-space register state and the stack pointer---a significant engineering challenge.
\end{principle}

\section{Exercises}

\begin{exercise}
	Explain why SMEP does not prevent ROP within the kernel. What does SMEP prevent that ROP cannot bypass?
\end{exercise}

\begin{exercise}
	A Linux device driver has a heap overflow in an \texttt{ioctl} handler. The kernel is compiled with SLAB\_FREELIST\_RANDOM and CONFIG\_SLUB. What additional information does an attacker need to exploit this?
\end{exercise}

\begin{exercise}
	Read about CVE-2022-0847 (Dirty Pipe). What was the vulnerability? What privilege did it grant? How did it bypass existing mitigations?
\end{exercise}

\begin{exercise}
	Design a seccomp filter for a process that only needs to read files (\texttt{read}, \texttt{openat}, \texttt{mmap}, \texttt{munmap}, \texttt{write} for stdout). Would this prevent exploitation of a kernel UAF bug triggered from this process? Why or why not?
\end{exercise}
